% \documentclass[11pt,respuestas,a4]{aleph-examen}
\documentclass[10pt,a5]{aleph-examen}

% -- Paquetes extra
\usepackage{aleph-comandos}
\usepackage{systeme}

% -- Datos
\institucion{Facultad de Ciencias Exactas, Naturales y Ambientales}
\carrera{Catálogo STEM}
\asignatura{Matemática III}
\tema{Recuperación -- Criterio 2.2}
\autor{Andrés Merino}
\fecha{Periodo 2026-1}

\logouno[0.16\textwidth]{Logos/logoPUCE_04_ac}
\definecolor{colortext}{HTML}{1957d1}
\definecolor{colordef}{HTML}{1957d1}
\fuente{montserrat}


\begin{document}

\encabezado

\vspace{-5mm}
% \section*{Indicaciones}
% \begin{itemize}[leftmargin=*]
% \item
%     En esta actividad se evalúa si el estudiante (Criterio 2.2) calcula derivadas de funciones de varias variables.
% \item
%     Los ejercicios deben ser resueltos a mano y entregarse en hojas de forma física al docente, sin ayuda de resúmenes o asistentes computacionales.
% \item
%     Se encuentra prohibido el uso de cualquier otra fuente de información durante todo el examen.
% \item
%     No está permitido tener ningún dispositivo electrónico encendido durante el examen, incluyendo celulares, tablets, computadoras, audífonos, relojes inteligentes, ni similares.
% \item
%     En caso de considerar que existe un error en la pregunta o que esta se encuentra mal planteada, se debe indicar cuál es el error y justificarlo.
% \item
%     Todas las soluciones deben estar correctamente redactadas y explicadas.
% \end{itemize}

\section*{Ejercicios}

\begin{preguntas}

%%%%%%%%%%%%%%%%%%%%%%%%%%%%%%%%%%%%%%%%%%
\item
Considere la función
\[
    f(x,y) = \begin{cases}
        \dfrac{x^2 y}{x^4 + y^2} & \text{si } (x,y)\neq(0,0),\\[6pt]
        0 & \text{si } (x,y)=(0,0).
    \end{cases}
\]
\begin{enumerate}[leftmargin=*]
    \item Calcule los límites iterados cuando $(x,y)\to(0,0)$.\puntaje{3}
    \item Calcule el límite por la trayectoria $y=2x$.\puntaje{2}
    \item Calcule el límite cuando $(x,y)\to(0,0)$ por la trayectoria $y=x^2$.\puntaje{2}
    \item ¿Es $f$ continua en $(0,0)$? Justifique a partir de los incisos anteriores.\puntaje{3}
\end{enumerate}


\begin{respuesta}
\begin{enumerate}
    \item Calculemos los límites iterados. Primero cuando $y\to 0$ y luego cuando $x\to 0$: para $x\neq 0$ fijo,
    \[
        \lim_{y\to 0}\frac{x^2 y}{x^4+y^2}
        =
        \frac{x^2\cdot 0}{x^4+0}
        =
        0,
    \]
    de modo que
    \[
        \lim_{x\to 0}\left(\lim_{y\to 0}\frac{x^2 y}{x^4+y^2}\right)
        =
        \lim_{x\to 0} 0
        =
        0.
    \]
    Ahora cuando $x\to 0$ y luego cuando $y\to 0$: para $y\neq 0$ fijo,
    \[
        \lim_{x\to 0}\frac{x^2 y}{x^4+y^2}
        =
        \frac{0\cdot y}{0+y^2}
        =
        0,
    \]
    de modo que
    \[
        \lim_{y\to 0}\left(\lim_{x\to 0}\frac{x^2 y}{x^4+y^2}\right)
        =
        \lim_{y\to 0} 0
        =
        0.
    \]

    \item Consideremos la trayectoria $y=2x$. Entonces:
    \[
        f(x,2x)
        =
        \frac{x^2\cdot 2x}{x^4+4x^2}
        =
        \frac{2x^3}{x^2(x^2+4)}
        =
        \frac{2x}{x^2+4}.
    \]
    Por tanto:
    \[
        \lim_{x\to 0} f(x,2x)
        =
        \lim_{x\to 0}\frac{2x}{x^2+4}
        =
        \frac{0}{4}
        =
        0.
    \]

    \item Consideremos la trayectoria $y=x^2$. Entonces:
    \[
        f\bigl(x,x^2\bigr)
        =
        \frac{x^2\cdot x^2}{x^4+\bigl(x^2\bigr)^2}
        =
        \frac{x^4}{x^4+x^4}
        =
        \frac{x^4}{2x^4}
        =
        \frac{1}{2}.
    \]
    Por tanto:
    \[
        \lim_{x\to 0} f\bigl(x,x^2\bigr)
        =
        \frac{1}{2}.
    \]

    \item Como los límites por distintas trayectorias son diferentes (por $y=2x$ se obtiene $0$, pero por $y=x^2$ se obtiene $\tfrac{1}{2}$), se concluye que
    \[
        \lim_{(x,y)\to(0,0)} f(x,y)
    \]
    no existe. Por lo tanto, $f$ no es continua en $(0,0)$.\qedhere
\end{enumerate}
\end{respuesta}


%%%%%%%%%%%%%%%%%%%%%%%%%%%%%%%%%%%%%%%%%%
\item
Sea $f(x,y)=y^2 e^x\cos(x)$.
\begin{enumerate}[leftmargin=*]
    \item Calcule $\dfrac{\partial f}{\partial x}$ y $\dfrac{\partial f}{\partial y}$.\puntaje{4}
    \item Calcule $\nabla f(P)$ en el punto $P=(0,2)$.\puntaje{2}
    \item Calcule la derivada direccional de $f$ en $P$ en la dirección del vector $\overrightarrow{PQ}$, donde $Q=(2,3)$.\puntaje{2}
    \item Calcule la derivada direccional máxima de $f$ en $P$ e indique la dirección en la que se alcanza.\puntaje{2}
\end{enumerate}


\begin{respuesta}
\begin{enumerate}
    \item Calculemos las derivadas parciales usando la regla del producto:
    \[
        \frac{\partial f}{\partial x}(x,y)
        =
        y^2\bigl(e^x\cos(x)+e^x(-\sen(x))\bigr)
        =
        y^2 e^x\bigl(\cos(x)-\sen(x)\bigr).
    \]
    \[
        \frac{\partial f}{\partial y}(x,y)
        =
        2y\,e^x\cos(x).
    \]

    \item Evaluamos en $P=(0,2)$:
    \[
        \frac{\partial f}{\partial x}(0,2)
        =
        4\cdot e^0\cdot(\cos 0-\sen 0)
        =
        4\cdot 1\cdot 1
        =
        4,
        \qquad
        \frac{\partial f}{\partial y}(0,2)
        =
        2(2)\cdot e^0\cdot\cos 0
        =
        4.
    \]
    Por lo tanto:
    \[
        \nabla f(P)
        =
        (4,\ 4).
    \]

    \item Calculamos el vector $\overrightarrow{PQ}$:
    \[
        u = \overrightarrow{PQ}
        =
        Q-P
        =
        (2,3)-(0,2)
        =
        (2,1).
    \]
    La derivada direccional de $f$ en $P$ en la dirección de $u$ es:
    \[
        f'(P;u)
        =
        \nabla f(P)\cdot u
        =
        (4,4)\cdot(2,1)
        =
        8+4
        =
        12.
    \]

    \item La máxima tasa de cambio de $f$ en $P$ se alcanza en la dirección del gradiente, es decir, tomando $u=\nabla f(P)=(4,4)$. La derivada direccional en esa dirección es:
    \[
        f'(P;u)
        =
        \nabla f(P)\cdot u
        =
        (4,4)\cdot(4,4)
        =
        16+16
        =
        32.\qedhere
    \]
\end{enumerate}
\end{respuesta}


%%%%%%%%%%%%%%%%%%%%%%%%%%%%%%%%%%%%%%%%%%
\item
Aplique la regla de la cadena para calcular las derivadas indicadas.
\begin{enumerate}[leftmargin=*]
    \item Calcule $\dfrac{dw}{dt}$ si
    $\quad
        w=x\cos(y),
        \quad
        x=t^2,
        \quad
        y=\ln(t).
    $\puntaje{3}

    \item Calcule $\dfrac{dz}{dt}$ si
    $
        \quad
        z=x^2y,
        \quad
        x=u+v^2,
        \quad
        y=uv,
        \quad
        u=e^t,
        \quad
        v=\ln(t).
    $\puntaje{5}
\end{enumerate}

\begin{respuesta}
\begin{enumerate}
    \item Aplicamos la regla de la cadena:
    \[
        \frac{dw}{dt}
        =
        \frac{\partial w}{\partial x}\cdot\frac{dx}{dt}
        +
        \frac{\partial w}{\partial y}\cdot\frac{dy}{dt}.
    \]
    Calculamos cada factor:
    \[
        \frac{\partial w}{\partial x}=\cos(y),\qquad
        \frac{\partial w}{\partial y}=-x\,\sen(y),\qquad
        \frac{dx}{dt}=2t,\qquad
        \frac{dy}{dt}=\frac{1}{t}.
    \]
    Por tanto:
    \[
        \frac{dw}{dt}
        =
        \cos(y)\cdot 2t + (-x\,\sen(y))\cdot\frac{1}{t}
        =
        2t\cos(y)-\frac{x}{t}\,\sen(y).
    \]

    \item Aplicamos la regla de la cadena considerando la dependencia completa:
    \[
        \frac{dz}{dt}
        =
        \frac{\partial z}{\partial x}\cdot\frac{dx}{dt}
        +
        \frac{\partial z}{\partial y}\cdot\frac{dy}{dt},
    \]
    donde a su vez:
    \[
        \frac{dx}{dt}
        =
        \frac{\partial x}{\partial u}\cdot\frac{du}{dt}
        +
        \frac{\partial x}{\partial v}\cdot\frac{dv}{dt},
        \qquad
        \frac{dy}{dt}
        =
        \frac{\partial y}{\partial u}\cdot\frac{du}{dt}
        +
        \frac{\partial y}{\partial v}\cdot\frac{dv}{dt}.
    \]
    Calculamos:
    \[
        \frac{\partial z}{\partial x}=2xy,\quad
        \frac{\partial z}{\partial y}=x^2,\quad
        \frac{\partial x}{\partial u}=1,\quad
        \frac{\partial x}{\partial v}=2v,\quad
        \frac{\partial y}{\partial u}=v,\quad
        \frac{\partial y}{\partial v}=u,
    \]
    \[
        \frac{du}{dt}=e^t,\qquad\frac{dv}{dt}=\frac{1}{t}.
    \]
    Entonces:
    \[
        \frac{dx}{dt}
        =
        1\cdot e^t+2v\cdot\frac{1}{t}
        =
        e^t+\frac{2v}{t},
    \]
    y
    \[
        \frac{dy}{dt}
        =
        v\cdot e^t+u\cdot\frac{1}{t}
        =
        ve^t+\frac{u}{t}.
    \]
    Por lo tanto:
    \[
        \frac{dz}{dt}
        =
        2xy\!\left(e^t+\frac{2v}{t}\right)+x^2\!\left(ve^t+\frac{u}{t}\right).\qedhere
    \]
\end{enumerate}
\end{respuesta}


\end{preguntas}

\end{document}
